\chapter{Bruises}

\section*{Definition and Overview}
A bruise is a collection of blood beneath the skin that forms when small blood vessels are damaged by a bump or impact and leak blood into the surrounding tissues. Medically it is termed a contusion, or an ecchymosis when it is large and flat, while a discrete, raised collection of blood is called a haematoma. Most bruises are minor and resolve without treatment, but they can also indicate a more serious underlying injury such as a fracture, a muscle tear, or internal bleeding, and unexplained or excessive bruising may signal a bleeding disorder. Recognising which bruises are benign and which need assessment is an important practical skill.

\section*{Epidemiology}
Bruising is one of the most common injuries of everyday life and affects people of all ages. Children bruise readily during normal play, and older adults bruise more easily because the skin becomes thinner and blood vessels more fragile with age. People taking anticoagulant or antiplatelet medicines bruise more easily, as do people with low platelet counts, heavy alcohol intake, or bleeding disorders. Bruises from sport, falls, and minor accidents account for a very large proportion of the minor injuries seen in primary care and emergency departments, and women tend to bruise more easily than men.

\section*{Aetiology and Pathophysiology}
A bruise forms when an impact crushes or stretches the small blood vessels beneath the skin sufficiently to make them leak. Blood escapes into the tissues, collects, and clots, producing a tender, discoloured swelling. Over subsequent days the body breaks down the blood pigments, which is why a bruise changes colour from red to purple, then blue, green, and yellow-brown before fading. Conditions that increase bruising include factors that weaken vessel walls, impair clotting, or reduce platelet numbers.

\begin{itemize}
    \item \textbf{Blood-thinning medicines:} anticoagulants such as warfarin and direct oral anticoagulants, and antiplatelet agents including aspirin
    \item \textbf{Bleeding disorders:} haemophilia, von Willebrand disease, and other inherited or acquired clotting defects
    \item \textbf{Thrombocytopenia:} a low platelet count, from disease or medications
    \item \textbf{Nutritional deficiency:} vitamin C or vitamin K deficiency
    \item \textbf{Ageing:} thinning of the skin and weakening of vessel walls
    \item \textbf{Corticosteroid use:} long-term use thins the skin and vessels
    \item \textbf{Alcohol:} heavy drinking impairs platelet function and clotting
\end{itemize}

\section*{Clinical Features}
The cardinal sign is a discoloured, tender area of skin that appears after an injury.

\begin{itemize}
    \item A red, purple, or blue mark that appears soon after a knock or impact
    \item Swelling and localised tenderness over the affected area
    \item A raised lump (haematoma) when a larger vessel has leaked
    \item Progressive colour change over one to two weeks as the bruise fades
    \item Pain on movement when the bruise overlies a joint or muscle
    \item Concerning features include bruising with severe pain, marked swelling, deformity, or inability to use the affected part
\end{itemize}

\section*{Differential Diagnosis}
Not every mark that resembles a bruise is a simple soft-tissue injury.

\begin{itemize}
    \item \textbf{Fracture or sprain:} an underlying bone or ligament injury hidden beneath the bruise
    \item \textbf{Muscle tear:} deep bleeding within a muscle causing swelling and loss of function
    \item \textbf{Bleeding disorders:} spontaneous or excessive bruising from clotting abnormalities or low platelets
    \item \textbf{Orbital fracture:} a black eye can indicate a fracture of the eye socket
    \item \textbf{Septicaemia or meningitis:} a spreading rash of small purple spots that does not fade under pressure
    \item \textbf{Purpura and petechiae:} pinpoint or larger bleeding spots from vasculitis, infection, or other disease
    \item \textbf{Non-accidental injury:} patterns of bruising inconsistent with the story, particularly in children or vulnerable adults
\end{itemize}

\section*{When to Seek Medical Care}
See a doctor if a bruise is very large or unusually painful, if several bruises appear without an obvious cause, if bruising is accompanied by nosebleeds or bleeding gums, or if you take blood thinners and sustain an injury. Any bruise that has not improved after about two weeks warrants review.

\textbf{Call 000 (or local emergency services) for:}
\begin{itemize}
    \item A bruise from a fall or impact with severe pain, deformity, or inability to move the limb
    \item Bruising after a head injury with confusion, vomiting, or drowsiness
    \item A rapidly spreading rash of small purple spots that does not fade when pressed
    \item Sudden severe swelling of a limb with pain and coldness
    \item Signs of internal bleeding, including dizziness, weakness, or fainting after an injury
\end{itemize}

\section*{Diagnosis and Investigations}
Most simple bruises are diagnosed from the history and examination alone; tests are reserved for severe or unexplained bruising.

\begin{itemize}
    \item \textbf{History:} the mechanism of injury, medications, alcohol intake, and any personal or family history of easy bleeding
    \item \textbf{Examination:} assessment of the size, tenderness, and any associated injury to joints, bones, or other structures
    \item \textbf{Blood tests:} platelet count and coagulation studies when easy bruising is reported
    \item \textbf{Imaging:} X-ray if a fracture is suspected, or ultrasound for a large or persistent haematoma
    \item \textbf{Specialist referral:} to a haematologist when an underlying bleeding disorder is suspected
\end{itemize}

\section*{Management}
Most bruises require no more than simple self-care, and treatment otherwise focuses on the underlying cause.

\begin{itemize}
    \item \textbf{Immediate first aid:} rest, ice wrapped in a cloth applied for 15--20 minutes at a time, and elevation of the injured part
    \item \textbf{Compression:} a firm bandage can limit swelling of a larger bruise
    \item \textbf{Pain relief:} paracetamol if needed; avoid aspirin because it can increase bleeding
    \item \textbf{Warmth:} gentle heat after the first two days may help the bruise fade more quickly
    \item \textbf{Observation:} monitoring for worsening pain, swelling, or loss of function
    \item \textbf{Treatment of the cause:} adjustment of blood-thinning medicines or investigation of an underlying bleeding problem under medical supervision
\end{itemize}

\section*{Complications}
Most bruises settle without difficulty, but complications occasionally arise.

\begin{itemize}
    \item A large haematoma that may require drainage
    \item Infection if the skin was broken at the time of injury
    \item Compartment syndrome, when bleeding within a limb muscle compresses nerves and blood supply
    \item A missed underlying injury, such as a fracture or ligament tear
    \item Myositis ossificans, a hardening of muscle after deep bruising, causing persistent swelling and pain
    \item Anaemia from severe, repeated bruising, particularly with bleeding disorders
\end{itemize}

\section*{Prognosis and Outcomes}
The vast majority of simple bruises heal fully within one to two weeks as the trapped blood is reabsorbed, with the colour changes gradually fading. Larger or deeper bruises can take several weeks to resolve and may leave a firm lump for a while. People who bruise easily because of medicines or an underlying condition generally do well once the cause is identified and managed. Persistent pain, swelling, or a residual lump after a bruise should be reassessed.

\section*{Prevention and Self-Care}
Bruising cannot always be prevented, but the risk of injury can be substantially reduced.

\begin{itemize}
    \item Use protective padding for sport and wear suitable, sturdy footwear on uneven ground
    \item Keep walkways clear and well lit, and use handrails on stairs
    \item Place non-slip mats in the bathroom and kitchen
    \item Check medicines and supplements that increase bleeding risk with a doctor or pharmacist
    \item Eat a balanced diet containing vitamin C and vitamin K
    \item Supervise children in play areas and remove obvious hazards
    \item Apply cold promptly after any knock to limit the size of the bruise
\end{itemize}

\section*{Sources and Further Reading}
The following organisations provide reliable information on bruises, first aid, and bleeding disorders for patients and healthcare students.

\begin{itemize}
    \item NHS. \textit{Bruises and bruises that won't heal: when to see a doctor.} https://www.nhs.uk/conditions/bruises/
    \item Mayo Clinic. \textit{Bruises: causes, symptoms, and home treatment.} https://www.mayoclinic.org/
    \item MedlinePlus. \textit{Bruises and contusions health information.} https://medlineplus.gov/bruises.html
    \item MSD Manual. \textit{Bruising: causes and evaluation in adults.} https://www.msdmanuals.com/
    \item Centers for Disease Control and Prevention. \textit{Injury prevention and first aid guidance.} https://www.cdc.gov/
    \item World Health Organization. \textit{Patient safety and injury prevention resources.} https://www.who.int/
\end{itemize}
